\chapter{Benchmarks, Baselines, and Metrics}
\label{ch:benchmarks}

The previous thirteen chapters built the machine and proved properties about it.
This chapter does something different: it lays out the \emph{instruments} that
measure whether the machine works. Three instruments, in order. The
\textbf{benchmarks} decide what we ask the system. The \textbf{baselines} decide
what we compare it against. The \textbf{metrics} decide how we score what comes
back. \Cref{ch:trajectory} and \Cref{ch:ablations} will report what these
instruments measured; this chapter is about the instruments themselves, because a
result is only as trustworthy as the ruler that produced it.

\begin{intuition}[Why a chapter on measurement]
It is easy to make a hallucination-reduction system look good by choosing a
favourable benchmark, a weak baseline, or a metric that rewards the behaviour the
system happens to produce. \caem{} commits in advance to a fixed panel of
benchmarks, a panel of eight baselines that each isolate one competing idea, and a
metric suite that scores honesty and accuracy separately. The discipline is the
point: every number in the next two chapters is produced by an instrument defined
here, on data the system never trained on, under a protocol that is identical for
\caem{} and for every baseline it is compared against.
\end{intuition}

\section{The benchmark panel}
\label{sec:bench-panel}

\caem{} is evaluated on a fixed panel of five question-answering benchmarks. They
are not interchangeable: each stresses a different failure surface, and they are
split into two roles that the whole continual-learning story depends on.

\begin{intuition}[Training benchmarks and transfer benchmarks]
A self-improving system has an obvious way to cheat: it can get better at exactly
the tasks it trains on while quietly getting worse everywhere else. To catch this,
the panel is split. Three benchmarks are \emph{training} benchmarks: the
self-improvement loop is allowed to learn from their answers. Two are
\emph{transfer} benchmarks: they are evaluated every cycle but the loop never
trains on them, so they measure whether an improvement generalises or is just
memorisation of the training distribution. A gain that shows up only on the
training benchmarks is suspect; a gain that also shows up on the held-out transfer
benchmarks is real.
\end{intuition}

\begin{table}[htbp]
\centering
\small
\begin{tabular}{@{}llll@{}}
\toprule
\textbf{Benchmark} & \textbf{What it tests} & \textbf{Answer surface} & \textbf{Role} \\
\midrule
FEVER~\cite{thorne-etal-2018-fever}                 & factual claim verification    & supports / refutes / not enough info & training \\
TriviaQA~\cite{joshi-etal-2017-triviaqa}            & open-text factoid recall      & free-text entity (many aliases)      & training \\
CommonsenseQA~\cite{talmor-etal-2019-commonsenseqa} & everyday commonsense          & five-option multiple choice          & training \\
TruthfulQA~\cite{lin-etal-2022-truthfulqa}          & common-misconception probing  & free-text answer                     & transfer \\
StrategyQA~\cite{geva2021strategyqa}                & implicit multi-step reasoning & yes / no                             & transfer \\
\bottomrule
\end{tabular}
\caption[The five-benchmark panel]{The evaluation panel. Three training benchmarks
feed the self-improvement loop; two transfer benchmarks are held out to measure
generalisation. All five are evaluated at every cycle.}
\label{tab:benchmarks}
\end{table}

\begin{precise}[What each benchmark stresses]
The five are chosen to span distinct surfaces so a single metric cannot be gamed.
\emph{FEVER} is a fact-checking benchmark with a bounded three-label space; the
answer is a verdict, not free text, so the failure mode is mislabelling.
\emph{TriviaQA} is open-text factoid recall with many acceptable aliases per
question, so the failure mode is fabrication of a plausible-but-wrong entity.
\emph{CommonsenseQA} is a five-option multiple-choice benchmark on everyday
reasoning with a bounded label space. \emph{TruthfulQA} is adversarial: its
questions are written so that the most fluent answer is a common human
misconception, so it specifically rewards a system that knows when not to commit.
\emph{StrategyQA} asks yes/no questions whose reasoning chain is implicit and must
be decomposed, stressing multi-step inference. The two transfer benchmarks
(TruthfulQA, StrategyQA) are exactly the ones the loop never sees during training.
\end{precise}

\begin{figure}[htbp]
\centering
\begin{tikzpicture}[
  font=\small\sffamily,
  train/.style={rectangle, rounded corners=4pt, draw=cbExample!70!black, fill=cbExample!12,
    line width=1pt, align=center, inner sep=4pt, text width=2.7cm, minimum height=1.0cm},
  trans/.style={rectangle, rounded corners=4pt, draw=cbScenario!70!black, fill=cbScenario!12,
    line width=1pt, align=center, inner sep=4pt, text width=2.7cm, minimum height=1.0cm},
  loop/.style={rectangle, rounded corners=6pt, draw=cbFormula!70!black, fill=cbFormula!8,
    line width=1.1pt, align=center, inner sep=5pt, text width=3.0cm, minimum height=1.0cm},
  flow/.style={-{Stealth[length=2.4mm]}, line width=1pt, draw=cbInk!75},
  held/.style={-{Stealth[length=2.4mm]}, line width=1pt, draw=cbScenario!75, dashed},
]
  \node[train] (fe) at (0,2.0) {FEVER\\{\footnotesize fact-checking}};
  \node[train] (tq) at (0,0.7) {TriviaQA\\{\footnotesize open-text factoid}};
  \node[train] (cq) at (0,-0.6) {CommonsenseQA\\{\footnotesize five-option MCQ}};
  \node[trans] (tf) at (0,-2.1) {TruthfulQA\\{\footnotesize misconception probe}};
  \node[trans] (sg) at (0,-3.4) {StrategyQA\\{\footnotesize implicit multi-step}};
  \node[loop] (sil) at (5.4,1.0) {self-improvement\\loop\\{\footnotesize trains on answers}};
  \node[loop, draw=cbInk!55, fill=cbInk!6] (ev) at (5.4,-2.0) {per-cycle\\evaluation\\{\footnotesize scores all five}};
  \draw[flow] (fe.east) -- (sil.west);
  \draw[flow] (tq.east) -- (sil.west);
  \draw[flow] (cq.east) -- (sil.west);
  \draw[held] (tf.east) to[bend right=8] (ev.west);
  \draw[held] (sg.east) to[bend right=12] (ev.west);
  \draw[flow] (fe.east) to[bend left=4] (ev.west);
  \draw[flow] (tq.east) to[bend left=2] (ev.west);
  \draw[flow] (cq.east) -- (ev.west);
  \node[font=\scriptsize, text=cbScenario!70!black, align=center] at (5.4,-3.6)
    {dashed: transfer benchmarks\\are never trained on};
\end{tikzpicture}
\caption[The training-versus-transfer split]{The panel split. The three training
benchmarks (green) feed the self-improvement loop and are also evaluated; the two
transfer benchmarks (orange) are evaluated every cycle but never enter training, so
they measure generalisation rather than memorisation. (Schematic.)}
\label{fig:panel}
\end{figure}

\subsection*{The panel was pruned for a principled reason}

The five-benchmark panel is what survived a cold-start screening. Three further
benchmarks were considered and removed before the main run, and the reason ties
directly back to the data-purity theorem of \Cref{ch:theory}.

\begin{bythenumbers}[Three benchmarks were removed, and why]
The purity theorem requires the verifier to clear a balanced-accuracy floor on the
claims a benchmark produces; below that floor, storing is worse than not storing. A
cold-start screening measured this and rejected three benchmarks. A multi-hop
benchmark's base correctness was about $0.09$, which demanded a verifier
discrimination far beyond what any signal achieved, so its episodes were
\emph{mathematically unstoreable}. An open-domain benchmark sat at about $0.16$ with
a verifier balanced accuracy below one half on every calibration fold, a
structural-failure benchmark. A four-option science benchmark was dropped as
redundant with the five-option commonsense benchmark already in the panel. All three
are kept on disk as registered exclusion evidence; none enters the active
trajectory. The lesson is that the panel is not a convenience sample, it is the set
of benchmarks on which the architecture's own precondition is met.
\end{bythenumbers}

\subsection*{Six disjoint pools, and why leakage would invalidate everything}

Within each benchmark, every sample is assigned to exactly one of six pools, and the
assignment is enforced to be leak-free. This is the single most important piece of
evaluation hygiene in the whole project.

\begin{precise}[The six pools]
Each benchmark's samples are partitioned into six disjoint pools: a \emph{seed} pool
that bootstraps the cold-start memory, a \emph{purity} pool that measures the
storage precision, a \emph{calibration} pool on which the composite of
\Cref{ch:composite} is fitted, per-cycle \emph{training} chunks that the
self-improvement loop consumes (one disjoint chunk per cycle), an \emph{evaluation}
pool that produces the per-cycle trajectory numbers, and a held-out \emph{test} pool
reserved for the final generalisation table. A sample that appears in training or
calibration must never reappear in evaluation or test, otherwise the reported
accuracy and calibration would be inflated by memorisation rather than measuring
generalisation. Transfer benchmarks populate only the evaluation and test pools;
their training, seed, calibration, and purity pools are empty by construction.
\end{precise}

\begin{fromcode}[Content-hash deduplication and the leakage guards]
Disjointness is guaranteed by a content hash, not by a dataset's native identifiers,
because source datasets reuse question text under drifting identifiers across
versions. Each sample's pool key is the first sixteen hexadecimal characters of a
SHA-256 digest of its question text, so two questions with identical text collide
into the same key regardless of how the dataset labelled them. Before any pool is
used, an assertion flattens all six pools and raises an error if a single content key
appears in more than one; a second assertion checks that no key appears in two
different benchmarks. Any duplicate halts the run immediately. The
deduplicate-then-slice order is what lets the evaluation and test numbers be read as
generalisation rather than recall.
\end{fromcode}

\begin{bythenumbers}[Pool sizes in the realised run]
In the realised trajectory each training benchmark allocated roughly one thousand
seed candidates, five hundred calibration samples, five hundred purity samples, a
per-cycle training chunk (two thousand for the two large benchmarks, seven hundred
for the smaller commonsense benchmark), three hundred evaluation samples, and up to
five hundred test samples. Each transfer benchmark allocated only its three hundred
evaluation samples and its test samples. The headline figure to carry is the
\emph{evaluation fold}: three hundred samples per benchmark across five benchmarks,
so every cycle's trajectory number is computed over fifteen hundred held-out
questions. The calibration fold that the composite is fitted on is a separate five
hundred per training benchmark, fifteen hundred in total, and is disjoint from the
evaluation fold.
\end{bythenumbers}

\subsection*{The retention probes}

One more measurement surface deserves its own mention, because it is what the
retention guard of \Cref{ch:cycle} reads. It is deliberately \emph{not} one of the
five panel benchmarks.

\begin{precise}[Three held-out probes, anchored once]
The retention guard scores the fine-tuned model on three probes that span different
abilities: a general-knowledge multiple-choice probe, a held-out open-text factoid
probe, and a held-out commonsense multiple-choice probe, two hundred questions each.
Their cold-start scores are anchored once, to about $0.61$, $0.395$, and $0.825$
respectively, and that anchor is reused as the denominator for every later cycle. A
cycle aborts if any probe's ratio against its anchor falls below the floor
$\rhomin = 0.93$. The probes are held out from both training and the panel
evaluation folds, so a high retention ratio cannot be bought by overfitting the
panel. This is the multi-modal guard that fired at the fourth cycle of
\Cref{ch:trajectory}, when the open-text factoid probe ratio fell to $0.886$.
\end{precise}

\section{The baselines: what we compare against}
\label{sec:bench-baselines}

A trajectory number in isolation means nothing. The question is always
\emph{compared to what}. \caem{} answers this with a panel of eight baselines, and
the panel is constructed so that each baseline isolates exactly one idea that a
sceptic might propose as the \emph{real} source of any gain.

\begin{intuition}[Each baseline removes one explanation]
Suppose \caem{} reports a gain over a plain language model. A reviewer can offer many
deflating explanations: maybe it is just the prompt, maybe just chain-of-thought,
maybe just retrieval, maybe just fine-tuning. The baseline panel is built to answer
each of these in turn. Every baseline runs on the \emph{same} frozen backbone as
\caem{}, so no comparison is contaminated by a stronger model; each adds exactly one
of the competing ingredients; and the gap between \caem{} and each baseline is the
value of everything \caem{} adds beyond that ingredient. No single baseline combines
all the ingredients, which is precisely the gap the architecture occupies.
\end{intuition}

\begin{table}[htbp]
\centering
\small
\begin{tabular}{@{}llp{6.6cm}@{}}
\toprule
 & \textbf{Baseline} & \textbf{What it isolates} \\
\midrule
B1 & zero-shot               & the empirical floor: the bare backbone with the same answer format, no reasoning trigger, no retrieval, no memory, no verifier \\
B2 & chain-of-thought        & the effect of a zero-shot reasoning trigger over the bare format~\cite{wei2022chain} \\
B3 & retrieval-augmented     & unconditional grounding: every query retrieves, no router, no verifier~\cite{karpukhin-etal-2020-dense} \\
B4 & reasoning $+$ retrieval & a reasoning trigger layered on top of retrieval \\
B5 & five-shot reasoning     & in-context demonstrations instead of fine-tuning~\cite{wei2022chain} \\
B6 & vanilla fine-tuning     & plain supervised fine-tuning with no retention guard, no memory, no verifier \\
B7 & active retrieval        & look-ahead-triggered retrieval during generation~\cite{jiang-etal-2023-active} \\
B8 & semantic-entropy        & sampling-uncertainty abstention from a single model family~\cite{farquhar2024semantic} \\
\bottomrule
\end{tabular}
\caption[The eight-baseline panel]{Each baseline runs on the same frozen backbone as
\caem{} and isolates one competing explanation for a gain. B1 is the floor; B3 and B7
isolate retrieval; B6 isolates training-time learning; B8 isolates sampling-based
abstention.}
\label{tab:baselines}
\end{table}

\begin{precise}[The matched protocol]
The comparison is fair only if every system is scored under identical conditions, so
three things are pinned. First, the \emph{backbone}: every baseline uses the same
frozen three-billion-parameter instruction model \caem{} uses, so a gain is never
attributable to a bigger model. Second, the \emph{evaluation fold}: baselines are
scored on the byte-identical, content-hash-matched evaluation samples \caem{} saw,
produced by the same seeded split. Third, the \emph{answer format}: every baseline
emits its answer through the same reasoning-then-answer scaffold \caem{} uses, so the
answer extractor reads the same shape from all systems. With these pinned, any
difference in score is attributable to mechanism, not to setup.
\end{precise}

\begin{pitfall}[The format floor: why ``raw'' exact match lies]
An earlier protocol let the generation-only baselines emit bare prose with no
explicit answer line. The answer extractor then found nothing to extract, and strict
exact match read close to zero \emph{even when the prose contained the correct
answer}. On the open-text factoid benchmark the bare zero-shot model was right in its
prose about fifty-nine percent of the time yet scored near zero on extracted exact
match. The fix is the matched protocol above: every baseline is required to commit to
an explicit answer line under the same instruction \caem{} follows, so its exact
match reflects what it knew rather than how it happened to format the output. The
book reports the matched-protocol numbers throughout, because a baseline made
artificially weak is as dishonest as a system made artificially strong.
\end{pitfall}

\begin{fromcode}[Scoring the baselines through the same verifier]
A baseline owns no verifier, so on its own it produces no hallucination signals and
no calibrated trust score. To place every system on the same honesty axis, each
baseline's question-and-answer pairs are replayed through the \emph{locked} \caem{}
verifier and composite, which recomputes the grounding, entailment, relevance, and
consistency signals and the four-outcome decision for that answer. The same locked
instrument scores \caem{} and every baseline, so any difference in the hallucination
metric is a property of the generator's output, not of a different ruler. This is the
standard practice for cross-system hallucination comparison: fix the judge, vary the
system under test.
\end{fromcode}

\begin{defence}[If an examiner asks: ``how much of the gain is just retrieval?'']
This is exactly what the panel is designed to answer, and the answer is built into
its structure. The contrast against the zero-shot floor (B1) isolates the joint
contribution of routing, verification, memory, and self-improvement together. The
contrast against unconditional retrieval (B3) isolates what the architecture adds
\emph{beyond} grounding every query. The contrast against vanilla fine-tuning (B6)
isolates what the deployment-time machinery adds beyond training-time learning. And
the difference between the B1 contrast and the B6 contrast decomposes the gain
between deployment-time mechanisms and training-time mechanisms. No reviewer can
attribute the result to retrieval alone, because B3 and B4 already hold retrieval
fixed; the remaining gap is everything else.
\end{defence}

\section{The metrics}
\label{sec:bench-metrics}

The final instrument is the scorer. \caem{} reports a small suite of metrics, and the
central design decision is that \emph{accuracy and honesty are scored on separate
axes}. A system can be more accurate, or more honest, or both, and conflating the two
is how the field has historically rewarded confident wrongness.

\subsection*{Exact match and token overlap}

\begin{precise}[Exact match, per task family]
Exact match is the headline accuracy axis and the basis for every paired comparison.
It is computed per task family after a shared normalisation that lowercases, strips
punctuation, removes the articles \emph{a}, \emph{an}, and \emph{the}, and collapses
whitespace, matching the standard convention. On the fact-checking benchmark it is
label-exact over the three verdicts; on the open-text factoid benchmark it is a match
against any of the gold aliases, with token-level overlap reported alongside for
partial credit; on the two multiple-choice and yes/no benchmarks it is label-exact.
The extractor reads the final answer out of the reasoning-then-answer scaffold before
scoring, so a correct answer that follows its reasoning is still counted.
\end{precise}

\begin{lstlisting}[caption={Exact match and answer extraction, condensed from \texttt{eval/metrics.py}.}, label={lst:em}]
def normalise(text):                       # SQuAD / TriviaQA convention
    text = text.lower().translate(PUNCT_TABLE)
    return " ".join(t for t in text.split() if t not in {"a", "an", "the"})

def extract_cot_answer(text):              # pull the answer out of the scaffold
    if "Answer:" in text:
        return text.split("Answer:")[-1].strip()
    m = search(r"the answer is[:\s]+(.+?)(?:[.!?]|$)", text)   # natural-language CoT
    return m.group(1).strip() if m else text.strip()

def exact_match(prediction, gold):
    return float(normalise(prediction) == normalise(gold))
\end{lstlisting}

\begin{pitfall}[The misconception benchmark needs a language-model judge]
On the misconception benchmark, surface-overlap scoring is actively misleading. The
benchmark provides reference answers, and a token-overlap score rewards an answer for
echoing the misconception's wording even when it asserts the false belief. A bare
model that fluently restates the common misconception can therefore score high on
overlap while being exactly wrong. For this benchmark the book uses a language-model
judge that reads the question and answer and rules on truthfulness, not token
overlap. The judged correctness is far lower, and far more honest, than the overlap
score: the same answers an overlap metric scores around $0.80$ are judged correct
only around $0.43$. Wherever the misconception benchmark's accuracy appears in this
book, it is the judged value, never the overlap value.
\end{pitfall}

\subsection*{The composite hallucination metric}

Accuracy says whether the answer was right. It says nothing about \emph{how} the
system failed when it was wrong. The composite hallucination metric is the honesty
axis, and it is a structured decomposition rather than a single rate.

\begin{precise}[Eight failure modes, equally weighted]
The composite hallucination metric is the equal-weighted mean over an eight-axis
failure-mode taxonomy. The eight axes are: confident confabulation (wrong while the
calibrated trust score is at least one half), factual fabrication (wrong while the
weakest atomic fact is ungrounded), logical fabrication (wrong while the reasoning
does not entail the answer), off-topic answering (wrong while the answer is
irrelevant to the question), defensive evasion (wrong while emitting a stock
``cannot determine'' phrase), template leakage (a placeholder string leaked from the
prompt), false refusal (the answer was actually correct yet the system abstained or
discarded it), and length-padded over-generation (wrong while the answer is
implausibly long). A ninth axis, factual contradiction, is part of the taxonomy but
excluded from the default denominator, because the deployed natural-language
inference judge produces no contradiction class and the corresponding rate is
structurally zero. The subtypes overlap deliberately, so the metric is a mean
prevalence across failure modes, not the fraction of samples with any failure.
\end{precise}

\begin{figure}[htbp]
\centering
\begin{tikzpicture}[
  font=\footnotesize\sffamily,
  ax/.style={rectangle, rounded corners=3pt, draw=cbPitfall!65!black, fill=cbPitfall!10,
    line width=0.8pt, align=center, inner sep=3pt, text width=2.7cm, minimum height=0.85cm},
  hub/.style={circle, draw=cbFormula!70!black, fill=cbFormula!10, line width=1.1pt,
    align=center, inner sep=2pt, text width=1.7cm, font=\scriptsize},
]
  \node[hub] (h) at (0,0) {CHM\\{\scriptsize mean of 8}};
  \node[ax] (a1) at (-4.6,2.3) {confident\\confabulation};
  \node[ax] (a2) at (-1.55,2.7) {factual\\fabrication};
  \node[ax] (a3) at (1.55,2.7) {logical\\fabrication};
  \node[ax] (a4) at (4.6,2.3) {off-topic\\answering};
  \node[ax] (a5) at (-4.6,-2.3) {defensive\\evasion};
  \node[ax] (a6) at (-1.55,-2.7) {template\\leakage};
  \node[ax] (a7) at (1.55,-2.7) {false\\refusal};
  \node[ax] (a8) at (4.6,-2.3) {length-padded\\over-generation};
  \foreach \a in {a1,a2,a3,a4,a5,a6,a7,a8}{\draw[cbInk!45, line width=0.7pt] (h) -- (\a);}
\end{tikzpicture}
\caption[The hallucination taxonomy]{The composite hallucination metric is the
equal-weighted mean of eight measured failure-mode rates. Each axis is a distinct way
to be wrong; weighting them equally avoids privileging one definition of
hallucination over another. A ninth axis (factual contradiction) is in the taxonomy
but is structurally zero under the deployed judge and excluded from the denominator.
(Schematic.)}
\label{fig:chm-axes}
\end{figure}

\begin{lstlisting}[caption={The composite hallucination metric, condensed from \texttt{eval/metrics.py}.}, label={lst:chm}]
SUBTYPES = ["confident_confabulation", "factual_fabrication",   # 8 measured axes
            "logical_fabrication", "off_topic", "defensive_evasion",
            "template_leak", "false_refusal", "over_long"]       # contradiction excluded

def chm(samples):
    rates = hallucination_subtypes(samples)        # one prevalence per axis
    return sum(rates[s] for s in SUBTYPES) / len(SUBTYPES)   # equal-weighted mean
\end{lstlisting}

\begin{workedexample}[A confident confabulation scores on one axis]
Take a wrong answer the system nonetheless stored, with a calibrated trust score of
$0.62$. It is incorrect and its trust is above one half, so it fires the confident
confabulation axis, the costliest failure: the model was wrong and apparently sure.
If its weakest atomic fact were also ungrounded it would additionally fire factual
fabrication, and a single bad answer can light up several axes at once. The metric
counts each axis's prevalence across the fold and averages the eight, so a system
that trades confident confabulations for honest abstentions lowers one axis without
inflating the others, and its composite falls. That is exactly the trade the gate of
\Cref{ch:composite} is built to make.
\end{workedexample}

\subsection*{The composite evaluation score}

The accuracy axis and the honesty axis can move in opposite directions, and so can
calibration and retention. A single headline scalar is useful only if it refuses to
let one strong axis hide a collapsed one.

\begin{precise}[A geometric mean of five axes]
The composite evaluation score is the geometric mean of five axes, each in the unit
interval with one as best: accuracy (pooled exact match), epistemic quality (one
minus the composite hallucination metric), retention (the worst of the three probe
ratios, capped at one), calibration (one minus twice the calibration error, floored
at zero), and verifier reliability (the verifier's balanced accuracy rescaled so
chance maps to zero). The mean is geometric, not arithmetic, on purpose: a near-zero
axis drags the whole score down, so a system cannot buy a good headline by excelling
on four axes while catastrophically forgetting on the fifth. Each axis is clamped
just above zero so a single uncomputed axis cannot zero the score outright. The score
is a visualisation and headline tool only; the book always reports the five component
axes beside it so the reader can reconstruct it.
\end{precise}

\begin{lstlisting}[caption={The composite evaluation score, condensed from \texttt{eval/metrics.py}.}, label={lst:ces}]
def ces_score(acc, epi, ret, cal, ver, eps=0.01):   # five orthogonal axes in [0,1]
    axes = [min(max(a, eps), 1.0) for a in (acc, epi, ret, cal, ver)]
    prod = 1.0
    for a in axes:
        prod *= a
    return prod ** (1.0 / len(axes))                 # geometric mean: any collapse hurts
\end{lstlisting}

\begin{pitfall}[Two conventions for the score, do not conflate them]
The five-axis score is fully defined only for a cyclic, calibrated, verifier-gated
system: retention, calibration, and verifier reliability are not properties a
single-pass baseline even has. So there are two conventions. The
\emph{system-internal} convention reports the full five-axis mean for \caem{} across
cycles, and is the one to read as the headline. A \emph{cross-system} convention
compares \caem{} and the baselines on only the two shared axes, accuracy and
epistemic quality, and defaults the other three to one for the baselines purely so
the arithmetic runs. Under that second convention a baseline can show a higher number
than \caem{}, but only because three of its axes were set to one by fiat rather than
measured. That is a bookkeeping artefact, never a finding; the baselines were not
measured on retention, calibration, or verifier reliability because they do not
implement them.
\end{pitfall}

\subsection*{Significance, and the rule for licensing a claim}

\begin{precise}[Paired tests, corrected across the panel]
Because every system is scored on the identical evaluation samples, the comparisons
are paired, and paired tests are used throughout. On the exact-match axis the test is
McNemar's test with a continuity correction, the maximum-likelihood test for matched
binary outcomes~\cite{mcnemar1947note}. On the hallucination axis, which is
continuous, the test is the paired Wilcoxon signed-rank test. Both report a
bias-corrected bootstrap ninety-five percent interval on the delta beside the
p-value, because a p-value alone does not convey effect size. Within each benchmark
the eight baseline contrasts are corrected for multiple comparisons by the Holm
procedure~\cite{holm1979simple}, which controls the family-wise error rate at higher
power than a plain Bonferroni correction.
\end{precise}

\begin{intuition}[When a claim is allowed to be called supported]
A single significant p-value is not enough to license an architectural claim, because
with enough samples a trivially small difference becomes significant. The book uses a
conjunction: a contrast counts as supporting only if its Holm-corrected p-value is
below five percent \emph{and} the bootstrap interval on the exact-match delta lies
entirely above two percentage points. Meeting one of the two is partial support;
meeting neither is unsupported. This pre-registered rule is what separates a result
worth building an architecture on from a statistical curiosity.
\end{intuition}

\subsection*{The decision distribution and the abstention lens}

\begin{precise}[Four outcomes, and what abstention buys]
Because \caem{} can decline to answer, accuracy alone misreads it. The four-outcome
decision distribution (store, defer, discard, abstain) is reported alongside
accuracy, and the per-outcome accuracy is strictly ordered: stored answers are the
most accurate, abstentions the least, with defer and discard between. The honest
reading of the misconception benchmark depends on this. A bare model is not licensed
to abstain, so it answers every adversarial question and pays the price in confident
wrongness; \caem{} abstains on roughly a quarter to a third of those questions. Its
strict exact match therefore looks lower, while its hallucination metric on the same
benchmark falls: the answer set is becoming more honest even as the strict-correctness
rate declines. Exact match is the wrong lens for honesty, which is precisely why the
book reports the honesty axis separately.
\end{precise}

\section{Reading the panel: what good looks like}
\label{sec:bench-reading}

The full trajectory is the subject of \Cref{ch:trajectory} and the ablations are the
subject of \Cref{ch:ablations}. But the instruments are easier to trust once we have
seen them produce a headline, so this section reads one real comparison end to end.

\begin{figure}[htbp]
\centering
\begin{tikzpicture}
\begin{axis}[
  name=emax, width=12.2cm, height=4.5cm,
  ybar, bar width=7pt,
  symbolic x coords={zero-shot, CoT, RAG, CoT-RAG, five-shot, vanilla-FT, active-ret, sem-entropy, CAEM-C5},
  xtick=data, x tick label style={font=\scriptsize, rotate=28, anchor=east},
  ymin=0.38, ymax=0.58, ylabel={pooled EM}, ytick={0.4,0.45,0.5,0.55},
  font=\small\sffamily, tick label style={font=\scriptsize},
  ymajorgrids, grid style={cbInk!8}, nodes near coords,
  nodes near coords style={font=\tiny, /pgf/number format/fixed, /pgf/number format/precision=3},
]
  \addplot[draw=cbInk!55, fill=cbInk!12] coordinates {
    (zero-shot,0.507) (CoT,0.500) (RAG,0.458) (CoT-RAG,0.464) (five-shot,0.405)
    (vanilla-FT,0.490) (active-ret,0.525) (sem-entropy,0.505) };
  \addplot[draw=cbExample!70!black, fill=cbExample!30] coordinates { (CAEM-C5,0.517) };
\end{axis}
\begin{axis}[
  name=chmax, width=12.2cm, height=4.5cm, at={(emax.south)}, yshift=-1.7cm, anchor=north,
  ybar, bar width=7pt,
  symbolic x coords={zero-shot, CoT, RAG, CoT-RAG, five-shot, vanilla-FT, active-ret, sem-entropy, CAEM-C5},
  xtick=data, x tick label style={font=\scriptsize, rotate=28, anchor=east},
  ymin=0.09, ymax=0.19, ylabel={pooled CHM (lower better)}, ytick={0.10,0.13,0.16,0.19},
  font=\small\sffamily, tick label style={font=\scriptsize},
  ymajorgrids, grid style={cbInk!8}, nodes near coords,
  nodes near coords style={font=\tiny, /pgf/number format/fixed, /pgf/number format/precision=3},
]
  \addplot[draw=cbPitfall!60!black, fill=cbPitfall!20] coordinates {
    (zero-shot,0.120) (CoT,0.123) (RAG,0.148) (CoT-RAG,0.142) (five-shot,0.136)
    (vanilla-FT,0.127) (active-ret,0.173) (sem-entropy,0.120) };
  \addplot[draw=cbExample!70!black, fill=cbExample!30] coordinates { (CAEM-C5,0.109) };
\end{axis}
\end{tikzpicture}
\caption[The deployed baseline panel]{The actual deployed comparison at the final
cycle. Top: pooled exact match across the five-benchmark fold; \caem{} sits in the
upper band, at parity with the strongest baseline (active retrieval edges it by under
one point). Bottom: pooled composite hallucination metric, where \caem{} is the
lowest (most honest) of every system, beating active retrieval by about thirty-seven
percent relative. The headline is read on both axes at once: accuracy parity with the
best baseline, and a clear honesty win over all of them.}
\label{fig:baseline-bars}
\end{figure}

\begin{bythenumbers}[The headline, on two axes]
At the final cycle \caem{} reaches a pooled exact match of $0.517$ and a pooled
hallucination metric of $0.109$ over the five-benchmark fold. Of the eight baselines,
the best on accuracy is active retrieval at $0.525$, edging \caem{} by under one
point, but its hallucination metric is $0.173$, the worst of any system. \caem{} is
the lowest on the hallucination axis of every system measured. At its best cycle the
accuracy reads $0.544$, ahead of all eight. The honest verdict is therefore precise:
accuracy at parity with the strongest baseline, and a decisive, statistically
significant advantage on confident wrongness, which is the quantity the architecture
was built to reduce.
\end{bythenumbers}

\subsection*{Decomposing the gain}

A single \caem{}-versus-floor number bundles two very different contributions: what
the architecture buys before any training, and what the self-improvement loop adds on
top. The book separates them with a three-way decomposition, because they answer
different questions.

\begin{figure}[htbp]
\centering
\begin{tikzpicture}
\begin{axis}[
  width=12.0cm, height=5.2cm, ybar, bar width=11pt,
  symbolic x coords={architecture, self-improvement, full},
  xtick=data, x tick label style={font=\small},
  ymin=-3, ymax=11, ylabel={pooled EM change (percentage points)},
  font=\small\sffamily, tick label style={font=\scriptsize},
  legend style={at={(0.97,0.97)}, anchor=north east, font=\scriptsize, draw=cbInk!25},
  ymajorgrids, grid style={cbInk!8},
  nodes near coords, nodes near coords style={font=\tiny, fill=white, inner sep=1pt, /pgf/number format/fixed, /pgf/number format/precision=1},
]
  \draw[cbInk!50, dashed] (axis cs:architecture,0) -- (axis cs:full,0);
  \addplot[draw=cbScenario!70!black, fill=cbScenario!22] coordinates {
    (architecture,3.4) (self-improvement,-0.9) (full,2.5) };
  \addplot[draw=cbExample!70!black, fill=cbExample!28] coordinates {
    (architecture,8.8) (self-improvement,-0.1) (full,8.7) };
  \legend{five-benchmark panel, three training benchmarks}
\end{axis}
\end{tikzpicture}
\caption[Where the accuracy gain comes from]{The actual three-way decomposition of
the exact-match gain over the zero-shot floor. ``Architecture'' is the cold-start
contrast (three tiers, verifier, and routing, before any fine-tuning);
``self-improvement'' is the slice the cyclic loop adds; ``full'' is the two combined.
On the five-benchmark panel almost all the accuracy gain is architectural and the
self-improvement slice is roughly flat, because the loop's value is on the honesty
axis instead, which \Cref{fig:chm-decomp} decomposes on the same three contrasts. On
the three training benchmarks the architectural gain is much larger (about nine
points).}
\label{fig:headline-decomp}
\end{figure}

\begin{figure}[htbp]
\centering
\begin{tikzpicture}
\begin{axis}[
  width=12.0cm, height=5.2cm, ybar, bar width=11pt,
  symbolic x coords={architecture, self-improvement, full},
  xtick=data, x tick label style={font=\small},
  ymin=-5, ymax=3, ylabel={pooled CHM change (percentage points)},
  font=\small\sffamily, tick label style={font=\scriptsize},
  legend style={at={(0.97,0.04)}, anchor=south east, font=\scriptsize, draw=cbInk!25},
  ymajorgrids, grid style={cbInk!8},
  nodes near coords, nodes near coords style={font=\tiny, fill=white, inner sep=1pt, /pgf/number format/fixed, /pgf/number format/precision=1},
]
  \draw[cbInk!50, dashed] (axis cs:architecture,0) -- (axis cs:full,0);
  \addplot[draw=cbScenario!70!black, fill=cbScenario!22] coordinates {
    (architecture,1.5) (self-improvement,-2.0) (full,-0.5) };
  \addplot[draw=cbExample!70!black, fill=cbExample!28] coordinates {
    (architecture,-1.9) (self-improvement,-1.5) (full,-3.4) };
  \legend{five-benchmark panel, three training benchmarks}
\end{axis}
\end{tikzpicture}
\caption[Where the honesty gain comes from]{The actual three-way decomposition of the
composite hallucination metric over the zero-shot floor, the honesty companion to
\Cref{fig:headline-decomp}; \textbf{lower is better, so a bar below the dashed line is
an improvement}. The contrasts are the same: architecture is the cold-start step, then
the self-improvement slice, then the two combined. The reading is the mirror image of
the accuracy figure. On the full five-benchmark panel the architecture step
\emph{raises} the hallucination metric slightly (a bar above the line), because
cold-start retrieval amplifies the adversarial misconception benchmark, and the
self-improvement step is what pulls it back down. On the three training benchmarks
both steps lower it and the combined reduction is about three and a half points (close
to a quarter relative). Where the accuracy figure showed the architecture doing the
work and the loop flat, this figure shows the loop doing the honesty work.}
\label{fig:chm-decomp}
\end{figure}

\begin{precise}[Architecture, self-improvement, and the honesty story]
The decomposition reads off three contrasts. The architecture contrast (the
zero-shot floor against cold-start \caem{}, before any training) lifts pooled exact
match by about three and a half points on the full panel and nearly nine points on
the training benchmarks: the three tiers, the verifier, and the router account for
most of the headline accuracy gain on their own. The self-improvement contrast
(cold-start \caem{} against the third cycle, architecture held fixed) is roughly flat
on exact match. This is not a disappointment, it is the expected result, because the
loop's job is honesty, not raw accuracy, and \Cref{fig:chm-decomp} shows this
directly. On the full panel the architecture step actually \emph{raises} the pooled
hallucination metric slightly, from $0.120$ to $0.135$, because cold-start retrieval
amplifies the adversarial misconception benchmark; the self-improvement step is what
pulls it back, from $0.135$ to $0.115$, about fifteen percent. On the three training
benchmarks both steps reduce it, and the combined drop is about a quarter. The loop
folds retrieval-found knowledge into the model and trims confident wrongness; it was
never
going to move strict exact match much, and the metric suite is designed to show that
distinction rather than hide it.
\end{precise}

\begin{defence}[If an examiner asks: ``so the loop barely helps accuracy?'']
Correct, and the book says so plainly rather than burying it. The self-improvement
loop is close to flat on pooled exact match, and that is the honest reading of the
decomposition. Its contribution is on the honesty axis: about a fifteen percent
reduction in the hallucination metric across the same cycles on the full panel, close
to a quarter on the training benchmarks, a growing share of
queries answered cheaply from the zero-shot tier as knowledge is distilled into the
weights, and a memory that survives every retention-guarded rollback. Reporting a
flat accuracy slice next to a real honesty gain is exactly the separation of axes
this chapter argued for. A metric suite that reported only accuracy would have called
the loop worthless; the suite the book uses shows what it actually bought.
\end{defence}

With the instruments defined, we can now read what they measured. \Cref{ch:trajectory}
walks the realised six-cycle trajectory benchmark by benchmark and cycle by cycle, and
\Cref{ch:ablations} removes one architectural piece at a time to price each component
against the panel and metrics established here.
